\section{Reduction of the remaining degree-twenty-two problem to abelian surfaces}
\label{sec:small-surface-reduction}

This section concerns the complex product $A=E^4$, with $E=\C/\mathcal O_K$ and the product polarization. It gives the complete geometric reduction used in W04 for the normalization $w=e_W[\Gamma_1]$: every proposed certificate of degree at most twenty-two can be replaced, without changing its class or degree, by sums of translates of abelian surfaces. The reduction includes the nonabelian surfaces and the exceptional abelian surfaces missed by a catalog of coordinate-pair graphs. W04 performs the resulting finite comparison.

\subsection{W01: small positive Hermitian matrices are sums of integral dyads}

An integral dyad is $vv^*$ with $v\in\mathcal O_K^n$; repetitions are allowed. Here $*$ means conjugate transpose.

\begin{lemma}[W01]
Every positive semidefinite integral Hermitian matrix over $\mathcal O_K=\Z[\omega]$ of size two and trace at most seven, or size three and trace at most seven, is a sum of integral dyads.
\end{lemma}

\begin{proof}
We first treat a binary matrix with diagonal entries $a\leq b$ and off-diagonal entry $x$. A zero diagonal forces its row to be zero. If $a=1$, subtract the dyad of $(1,\bar x)$; the remainder is the nonnegative integral diagonal entry $b-\Nm(x)$.

The case $a=2$ works for any $b$. Since $\mathcal O_K/(2)=\F_4$, either $x$ is divisible by two, or its residue is represented by an Eisenstein unit. Choose $r=0$ in the first case and a unit with $r\equiv\bar x\pmod2$ in the second. Put
\[
u=(\bar x+r)/2,\qquad v=(\bar x-r)/2.
\]
Then $u,v\in\mathcal O_K$, $u+v=\bar x$, and
\[
\Nm(u)+\Nm(v)=\frac{\Nm(x)+\Nm(r)}2\leq b.
\]
The inequality follows from $\Nm(x)\leq2b$; in the nondivisible case $\Nm(x)$ is odd. The two dyads of $(1,u)$ and $(1,v)$ give the prescribed first diagonal and off-diagonal entries. Add coordinate dyads to supply the remaining second diagonal.

Only $(a,b)=(3,3)$ or $(3,4)$ remains. The norms at most twelve are
\[
0,1,3,4,7,9,12.
\]
This follows directly from $\Nm(c+d\omega)=c^2-cd+d^2$. Every element of norm at most nine is a sum of at most three units: up to units and conjugation the nontrivial cases are $1-\omega$, $2$, $2-\omega$, and $3$. Sum the corresponding dyads $(1,\bar\epsilon)$ and fill the two diagonals with coordinate dyads. Norm twelve can occur only for $(a,b)=(3,4)$; its elements are associates of $2(1-\omega)$, and the matrix is already a dyad with coordinate norms three and four. This proves the binary assertion.

Now let $H$ be ternary. A zero diagonal reduces to the binary case. If a diagonal is one, subtract its integral dyad as above. The positive Schur complement is binary, integral, and has trace at most six. If every diagonal is at least two, permute coordinates to obtain
\[
H=\begin{pmatrix}
a&x&y\\ \bar x&2&z\\ \bar y&\bar z&2
\end{pmatrix},\qquad a\in\{2,3\}.
\]
The norm of $z$ is $0,1,3$, or $4$. Integral unitary diagonal changes of coordinates preserve the desired conclusion.

If $\Nm(z)=4$, normalize $z=2$. Positivity forces the second and third columns to agree, since their difference has zero norm. Thus $x=y$ and the problem is the pullback of the binary matrix with diagonal $(a,2)$, already treated.

If $\Nm(z)=3$, the lower binary block is positive unimodular. Normalize $z=1-\omega$. The vector $(1,-1)$ has norm one; extending it by $(0,1)$ and subtracting its integral orthogonal projection gives an integral orthogonal basis with both norms one. Therefore an integral congruence makes the lower block $I_2$. The resulting matrix is the sum of the dyads $(x',1,0)$ and $(y',0,1)$ and $a-\Nm(x')-\Nm(y')$ coordinate dyads. Positivity makes the latter integer nonnegative. Undo the integral congruence.

Suppose $\Nm(z)=1$, and normalize $z=1$. Put
\[
s=\omega-\omega^2,\qquad
F=\frac{2\Nm(x)+2\Nm(y)-\Tr(x\bar y)}3.
\]
The Schur complement says $a\geq F$. If $s\mid x+y$, the elements
\[
A=(y+\omega x)/s,\qquad B=x-A
\]
are integral. The two dyads of
\[
(A,1,-\omega),\qquad(B,1,-\omega^2)
\]
have lower block $\left(\begin{smallmatrix}2&1\\1&2\end{smallmatrix}\right)$, off-diagonal entries $x,y$, and first diagonal $F$. Add $a-F$ coordinate dyads.

If $s\nmid x+y$, its residue modulo three is a unit. The six Eisenstein units map bijectively to the six units of $\mathcal O_K/(3)$: they are distinct modulo three, since a nonzero difference of units has norm at most four. Choose a unit $r\equiv x+y\pmod3$ and put $t=(x+y-r)/3$. The three dyads of
\[
(t,1,1),\qquad(x-t,1,0),\qquad(y-t,0,1)
\]
give the desired lower block and off-diagonal entries, with first diagonal $F+1/3$. Indeed
\[
\Nm(t)+\Nm(x-t)+\Nm(y-t)=F+\Nm(r)/3.
\]
Here $\Nm(x+y)\equiv1\pmod3$, so $F\equiv2/3\pmod\Z$. Since $a$ is integral and $a\geq F$, one has $a\geq F+1/3$. Fill the remaining first diagonal.

Finally let $z=0$. The two binary constructions use first diagonal entries $\lceil\Nm(x)/2\rceil$ and $\lceil\Nm(y)/2\rceil$. Positivity gives $\Nm(x)+\Nm(y)\leq2a$. These constructions therefore suffice unless both norms are odd and their sum is $2a$. The only exceptions are the norm pair $(1,3)$, up to order, for $a=2$, and $(3,3)$ for $a=3$. Up to independent unit changes in the last two coordinates, the first exception is the sum of the dyads
\[
(1,1,1),\qquad(-\omega,-1,1),
\]
whose off-diagonal norms are one and three and whose last two coordinates are orthogonal. The second is the sum of the dyads
\[
(1-\omega,1,1),\qquad(0,1,-1).
\]
Independent unit changes give every required pair of phases. This exhausts the four cases and proves the lemma.
\end{proof}

\paragraph{The cycle-class interpretation.}
Associate a divisor class on $E^n$ with its Rosati-fixed matrix using the product principal polarization; see \cite[Proposition 17.2, p.~35]{MilneAV}. The matrix is integral because this polarization is an isomorphism, and the product Rosati involution is conjugate transpose. An effective divisor is nef: translate it to meet any fixed curve properly and use translation invariance of its numerical class. Its matrix $H$ therefore satisfies $v^*Hv\geq0$ for every $v\in\mathcal O_K^n$, by pullback along $E\to E^n$. Clearing denominators gives the same inequality on $K^n$, which is dense in $\C^n$, proving positive semidefiniteness. A dyad is represented by the pullback of the origin under the homomorphism with the corresponding conjugate Eisenstein coefficients; its fibers are sums of translates of abelian subvarieties, with their actual integral multiplicities.

For a curve cycle $C$ on $E^n$, use instead its Hermitian moment matrix. Its diagonal entries are the degrees of the maps to the elliptic coordinates, and its $(i,j)$ entry is the integral cross-correspondence coefficient obtained by pushing $C$ to $E_i\times E_j$. The K\"unneth decomposition and the principal self-duality of $E$ identify this coefficient with an element of $\End(E)=\mathcal O_K$. Analytically the matrix integrates the products of the pulled-back holomorphic one-forms, so it is positive semidefinite. These entries determine the full curve class: its Hodge type is $(1,1)$ in homology, and the Hermitian products span the dual space of forms of that type. The dyad $vv^*$ is the class of the effective pushforward of $E$ under the map with coordinate vector $v$.

The word ``pushforward'' is essential. If that map has degree $e$ onto its image, its cycle is $e$ times the reduced elliptic image. Thus a dyad always gives an integral effective cycle, including when the entries have a common factor. No primitive-vector assumption or division of a cycle coefficient is made. For a curve the product-polarization degree is the matrix trace; for a divisor on $E^3$ its surface degree is twice the trace. Consequently every replacement supplied by W01 preserves the exact degree and the homological class.

\subsection{W02: every relevant small surface class has an abelian replacement}

\begin{proposition}[W02]
Every irreducible surface in $E^4$ of degree at most fourteen has the class of an effective integral sum of translates of abelian surfaces of the same degree. The same conclusion holds for every irreducible surface of degree at most twenty-two with nonzero Weil projection.
\end{proposition}

\begin{proof}
Let $a_{ij}$ be its six coordinate-pair projection degrees. Then
\[
\deg_\theta V=2\sum_{i<j}a_{ij}.
\]
If one pair degree is one, the extension and rigidity argument of V02 identifies $V$ itself as a translated abelian graph. We may therefore assume no pair has degree one.

First suppose the degree is at most fourteen. If a coordinate is constant, $V$ is an effective divisor on a translate of $E^3$. Its Hermitian matrix has trace at most seven, so W01 applies. If every coordinate is nonconstant, algebraic dependence of coordinate function subfields is an equivalence relation, as in V05. There are at most three dominant pairs, since every positive pair degree is at least two. The only possible partition is $(1,3)$. Thus $V=E_i\times C$ for an integral curve in the other three coordinates. Its curve degree is at most seven. Apply the curve version of W01 and take the product with $E_i$.

Now assume $\deg_\theta V\leq22$ and its Weil projection is nonzero. Every coordinate is nonconstant, and there are at most five dominant pairs. The possible partitions are $(1,3)$, $(2,2)$, and $(2,1,1)$. The first gives zero Weil projection, since three of the four degree-one factors pull back from one curve.

For $(2,2)$ the surface is $C\times D$ in the two coordinate pairs. Both curve degrees are at least two, their product is at most eleven, and each is therefore at most five. The binary part of W01 replaces each curve class by an effective sum of elliptic curve classes. Their products give the desired surface replacement. If one dependent pair is the first or second fixed $E^2$ factor, the Weil projection is already zero; this does not cause an additional case.

For $(2,1,1)$ the dependent pair must be a cross pair, since the wedge of two holomorphic forms from a curve vanishes. Relabel it $(1,3)$, and let $C$ be its curve image. Write $n_1,n_3$ for its two projection degrees. For the generic fiber over the normalization of $C$, write $b_2,b_4$ for the degrees of its two elliptic projections, including multiplicities. Then
\[
a_{12}=n_1b_2,\quad a_{14}=n_1b_4,\quad
a_{23}=n_3b_2,\quad a_{34}=n_3b_4.
\]
The remaining dominant pair has degree $a_{24}\geq2$. The four displayed integers are at least two and sum to at most nine. Since they form a rank-one $2\times2$ array, they cannot consist of three twos and one three. They are all two. Consequently
\[
n_1=n_3=n,\qquad b_2=b_4=b,\qquad nb=2,
\qquad a_{24}=d\in\{2,3\}.
\]

If $(n,b)=(1,2)$, the curve $C$ is a translated unit graph. The surface is an effective divisor on $C\times E_2\times E_4\simeq E^3$. Its Hermitian diagonal is $(d,2,2)$. W01 decomposes its divisor class into integral dyads; the corresponding abelian divisors, embedded in $E^4$, give the replacement. No degree-one projection to $C\times E$ is assumed here: both those projection degrees are two.

If $(n,b)=(2,1)$, the geometric generic fiber is an integral unit graph in $E_2\times E_4$. Indeed its conjugate geometric components have equal projection degrees; total degree one in both projections permits only one component, of multiplicity one. A degree-one projection identifies its normalization with $E$, and the other degree-one map is a translated automorphism. Its graph direction is a fixed unit $u\in\mathcal O_K^\times$, since the endomorphism ring is constant. Quotient by this fixed graph direction. The quotient map is split, for it uses $x_4-u x_2$ in the last two coordinates. The image of $V$ is an integral curve $C'\subset E^3$, and
\[
V=q^{-1}(C')\simeq C'\times E
\]
under this splitting: the equality holds first over the generic point of $C$ and then by closure. The three projection degrees of $C'$ are $(2,2,d)$; the last equals $a_{24}$ because the fibers of the split quotient have degree one. Apply the curve version of W01 to $C'$, and pull back its effective elliptic-cycle replacement by $q$.

Every replacement is an equality of homological classes by the matrix interpretation. It is effective with integral multiplicities, and equality of classes preserves its ambient $\theta$-degree. This proves both assertions.
\end{proof}

For the remaining $w$ problem, a certificate with maximum cost below twenty-four has multiplier three. Indeed $3$ divides the multiplier, and the complex calibration $|z(Z)|\leq\deg_\theta Z/8$ forces $D\geq4N$ for a representation of $Nw$. Both numerators must have nonzero Weil component, since an individual degree-at-most-twenty-two cycle has $|z|<3$. Thus every zero-Weil component has degree at most fourteen, while every nonzero-Weil component has degree at most twenty-two. W02 replaces all of them. This proves the claimed reduction of the full remaining problem to effective sums of abelian surfaces.

\subsection{W03: the abelian surfaces missed by a graph-only catalog}

\begin{proposition}[W03]
A nonzero-Weil abelian surface of degree at most twenty-two is a graph, or transposed graph, of an isogeny $h:E^2\to E^2$ with
\[
\deg h\in\{1,3,4\},\qquad
\sum_{i,j}\Nm(h_{ij})\leq10-\deg h.
\]
An abelian surface of degree at most fourteen either has a coordinate-pair projection of degree one or has degree fourteen and, up to permutation of coordinates, is
\[
E\times\{(s x,2u x):x\in E\}\times\{0\},
\qquad \Nm(s)=3,\quad u\in\mathcal O_K^\times.
\]
Translations do not change these class assertions. The latter surfaces have zero Weil projection and must be included in a complete catalog.
\end{proposition}

\begin{proof}
Every complex subtorus of $E^4$ is stable under scalar $\omega$: its tangent subspace is complex linear and its intersection lattice is stable under $\mathcal O_K$. The surface lattice is therefore a free rank-two $\mathcal O_K$ module. Choosing a basis identifies the surface with $E^2$, and its coordinate maps with Eisenstein matrices. Pair-projection degrees are norms of the $2\times2$ minors. Nonzero Weil projection forces both projections to the fixed two $E^2$ factors to be dominant. If both degrees were at least three, the Hodge-index bound would give
\[
D=2(a+b+c)\geq2(a+b+2\sqrt{ab})\geq24.
\]
An Eisenstein norm cannot equal two, so one projection degree is one. The graph degree formula gives $D=2(1+d+T)$ and $T\geq2\sqrt d$. The possible positive Eisenstein norms are then $d=1,3,4$, with the stated budget.

For the second assertion, the surface lattice is saturated in $\mathcal O_K^4$. Its six Pl\"ucker minors consequently generate the unit ideal. If none is a unit, there are at most two nonzero minors, since each positive norm is at least three and their norm sum is at most seven. One nonzero minor cannot be primitive. Two minors of norm three have a common factor above three. Thus exactly two are nonzero, of norms three and four, and the degree is fourteen. The Pl\"ucker relation forces these minors to share a coordinate: two disjoint nonzero minors would leave an uncancelled nonzero product. Hence their primitive exterior vector has the form
\[
e_i\wedge(s e_j+2u e_k)
\]
after multiplication by a unit and coordinate permutation. The two coefficients are coprime, so this is the asserted embedded product. Its fourth coordinate is constant, giving zero Weil projection. Conversely these products have pair degrees three and four and thus degree fourteen, verifying the exception exactly.
\end{proof}

As a computational check of W01, an independent exact enumeration found $1171$ positive semidefinite matrices with diagonal $(2,2,2)$ and $1963$ with diagonal $(3,2,2)$, and produced integral dyad decompositions for all of them. The proof above does not depend on accepting that enumeration. The script and output are retained as a reproducibility check, not as a formal proof certificate.
